\section{Interference software must preserve the population and the policy}
\label{sec:exec24-interference-reuse}

A bank considering outreach to connected borrowers must decide whose outcome it
wants to improve and which assignments it can change. A household contact, a
branch allocation and a directed financial relationship define different
interventions. MossSpider 0.0.3 and inferference 1.0.3 provide useful, but
different, implementation references. The former exposes a network-based
targeted procedure; the latter constructs grouped allocation inverse-probability
weights. Their selected commits are respectively
\texttt{62e6bd20364c7ef01b6964fd79f21be84e5ac579} and
\texttt{b0c785e3f60541b993ff4e7324672b761a5f9723}. The comparison below is
an inspection of authentic source, not a native benchmark or an assertion that
their unrestricted statistical ambitions have been qualified.
\citep{exec24_interference_mossspider,exec24_interference_inferference}

\subsection{A network population is not automatically a population of groups}
Let $A_i$ be a binary contact decision and let $G_{ij}=1$ mean that borrower
$j$'s treatment contributes to the exposure summary for borrower $i$. Define
$S_i(A)=\sum_jG_{ij}A_j$. A proposed network policy specifies a law for the
whole assignment vector and a population $E$ of eligible outcomes. A candidate
target is
\[
 \psi_E(\pi)=\frac{1}{|E|}\sum_{i\in E}E_\pi[Y_i(A_i,S_i(A))].
\]
This notation states a target; it does not identify it. Consistency, the
adequacy of the exposure mapping, measured confounding, support and a justified
sampling or network law remain substantive requirements. An API accepting an
adjacency matrix supplies none of them.

The selected MossSpider constructor requires integer node identifiers, refuses
self-loops and nonbinary treatment, saves an original-identifier mapping and
constructs an unweighted adjacency matrix. Its exposure helper multiplies rows
of that matrix by the ordered treatment vector. For the directed chain
$0\to1\to2$ and treatment vector $(1,0,0)$, row-neighbor exposure is
$(0,0,0)$; using incoming neighbors gives $(0,1,0)$. Both are calculable, but
they answer different banking questions. An adapter must declare the meaning
of an arrow, preserve node ordering and test an asymmetric graph. Automatically
symmetrizing a relationship network changes the intervention.
\citep{exec24_interference_mossspider,exec24_interference_reference}

For grouped partial interference, let group $g$ contain $n_g$ people and let
$\mu_g(\alpha)$ be its average outcome under allocation $\alpha$. Two
reasonable targets are
\[
 \psi_{\mathrm{group}}=\frac1M\sum_{g=1}^M\mu_g(\alpha),
 \qquad
 \psi_{\mathrm{person}}=\frac{\sum_g n_g\mu_g(\alpha)}{\sum_g n_g}.
\]
The inspected inferference point-estimate path first constructs group means
and then averages their weighted contributions with denominator $M$. Consider
two groups of sizes two and eight with stipulated outcome means USD~100 and
USD~0. Equal-group averaging gives USD~50, whereas equal-person averaging gives
USD~20. Neither is a software error. Substituting one for the other would be a
product error if a bank's objective were total borrower benefit. The treatment-
specific group contribution also divides by total group size, rather than the
number treated; its later allocation normalization must remain part of the
same estimand. An integration cannot compare the two packages simply because
both return a field called an effect.
\citep{exec24_interference_inferference,exec24_interference_reference}

\subsection{Eligibility restrictions change the intervention itself}
MossSpider's degree restriction selects the outcome population and also affects
policy simulation. The inspected generator keeps excluded nodes at their
observed treatment while resampling eligible nodes. This is a mixed policy,
not an all-network stochastic intervention. In a star network, suppose the hub
is excluded by a degree threshold, its observed treatment is one, and each
eligible leaf has stipulated outcome $10A_{\mathrm{hub}}$. Keeping that hub
fixed gives ten per leaf. Resampling the hub with probability $1/4$ gives
$5/2$. The distinction persists even if only the leaves' outcomes are averaged.
Eligibility, the treatment of excluded neighbors and the resulting target
population must therefore travel together in the request and response.
\citep{exec24_interference_mossspider,exec24_interference_reference}

Boundary behavior also requires an explicit contract. The selected fit
documentation describes a one-sample shortcut for deterministic policies, but
the selected body rejects uniform all-zero and all-one plans. It rejects an
integer argument even when its mathematical value could be a probability. A
wrapper should expose these selected-version restrictions, rather than promise
the documented shortcut or silently perturb zero into a small positive number.
The source includes tests for several input and model refusals; this inspection
does not execute those tests or claim that every boundary is covered.
\citep{exec24_interference_mossspider}

\subsection{Allocation mass must appear exactly once}
For a group assignment $a$, denote its observed probability by $p_\theta(a)$
and its proposed policy probability by $p_\alpha(a)$. The desired importance
ratio is $w_\theta(a;\alpha)=p_\alpha(a)/p_\theta(a)$ when both probabilities
refer to the same assignment and conditioning information. The inspected
default inferference integrand already divides the observed Bernoulli factors
by the policy factors. Consequently its integrated value represents a
probability ratio under that default construction. The reciprocal branch and
the branch that multiplies by policy mass are not interchangeable.
\citep{exec24_interference_inferference}

An exact no-random-effect witness makes the normalization visible. For two
treated members with common observed probability $q=1/2$ and proposed
probability $\alpha=1/4$, observed mass is $1/4$ and policy mass is $1/16$.
The ratio integral is four, so its reciprocal is $1/4$. Multiplying by policy
mass again gives $1/64$. If $q$ is parameterized by a common logit coefficient
$\beta$, then
\[
 \frac{\partial w}{\partial\beta}=-2(1-q)w,
\]
which equals $-1/4$ in the first construction and $-1/64$ in the duplicated-mass
construction. This is an authored algebraic example of two different
normalizations, not an observation of an executed package failure.
\citep{exec24_interference_reference}

The implementation review must follow the entire argument chain. The selected
derivative-array helper passes a false allocation-integration flag to its
derivative helper. That helper consumes the formal argument but does not place
it in the constructed gradient arguments; the downstream weight function then
uses its default. Reading only the first call would therefore misdescribe the
executed semantics. In a future native qualification, both flag values, the
default and a custom integrand must be checked against a known integral and an
independently calculated derivative. The selected repository includes numerical
weight and derivative tests, but their existence alone does not settle every
flag/default combination. A product should record the resolved settings and
refuse a score/weight mismatch.
\citep{exec24_interference_inferference}

\subsection{Numerical output does not settle uncertainty or missingness}
The grouped weight calculation catches an integration error and returns a
missing value. Later group aggregation uses a sum with missing values removed
while retaining the original group count. If three stipulated group
contributions are $10$, missing and $30$, this expression yields $40/3$. The
complete-case mean is $20$, and a hypothetical complete contribution of $60$
for the missing group would give $100/3$. None is a justified replacement for
the missing scientific contribution without further assumptions. The wrapper
must report submitted and usable groups, retain the failure and qualify or
refuse the target; it must not reinterpret a finite result as complete data.
\citep{exec24_interference_inferference,exec24_interference_reference}

The inspected grouped robust calculation also solves a score crossproduct
matrix, and the returned intervals use normal quantiles. A matrix inverse
requires numerical diagnostics; a label such as robust requires a justified
group-sampling law. In the network implementation, one variance uses squared
weighted residuals and another uses their adjacency quadratic form with the
diagonal filled. A symmetric path with residual vector $(1,-2,1)$ gives
residual-square expression $2/3$ and adjacency expression $-2/9$. This shows
that the bare finite-sample quadratic need not be positive semidefinite. It
does not show that these residuals arose from the fitted procedure or refute
any stated population theorem. Negative or undefined variance is nevertheless
a necessary refusal/diagnostic fixture, alongside simulation under a separately
justified network law. Clipping it to zero would conceal the issue.
\citep{exec24_interference_mossspider,exec24_interference_inferference,exec24_interference_reference}

\subsection{Separate reuse decisions and a bounded public prototype}
The proposal retains both packages as references. An original-language worker
is a conditional prototype route: it avoids an immediate scientific rewrite,
but still needs exact identity, default, missingness, policy and uncertainty
conformance. An exact finite-policy oracle supplies a useful baseline. A
whole-library port is not justified by the present evidence. It would add
maintenance and scientific obligations before demonstrating an advantage.
The two reference choices do not approve either as a general contagion or
dynamic banking network engine.
\citep{exec24_interference_reference}

MossSpider's inspected MIT notice requires notice preservation and disclaims
warranty; its selected setup lists several dependencies without fixed versions.
inferference's DESCRIPTION declares GPL version two or later and identifies its
R dependencies and maintainer. These are source facts, not a complete
distribution-rights determination. A platform owner must review the actual
installed distribution and transitive dependencies before shipping. Selected
commit and publication dates establish historical identity, not future support.
\citep{exec24_interference_mossspider,exec24_interference_inferference}

The first public prototype should prespecify asymmetric direction, unequal
group size, excluded-hub, deterministic-boundary, known-integral, derivative,
missing-group and singular/negative-variance fixtures. Scientific owners must
define the target and sampling law; adapter owners preserve identities and
failure states; evaluation owners separate native conformance from uncertainty
coverage. Every attempt contributes to cost. For a stipulated USD~2,400 total,
100 submitted requests and 60 usable outputs imply USD~24 per submission and
USD~40 per usable output. These figures illustrate denominators, not vendor
prices or bank returns. A later expensive bank experiment needs a qualified
public prototype and a credible economic case; neither follows merely from
the availability of these two interfaces.
\citep{exec24_interference_reference}
